% AI-produced by GPT-6 Astra (OpenAI, ChatGPT Pro mode), 19 November 2026.
% Parent manuscript supplies amsmath, amssymb, amsthm, mathtools, natbib.
\section{D. Real fibre functors, slope parity and the involution}
\label{sec:real-weil-D}

\paragraph{Conventions and scope.}
We use geometric Frobenius on cohomology, so a weight-$w$ Weil number has absolute
value $q^{w/2}$ and the positive Tate character is the character of $\mathbb Q(-1)$,
with Frobenius $q$. We write character groups additively. For a torus $T$ over
$\mathbb Q$ of the kind considered on 15:117--125, put
\[
 M=X^*(T_{\overline{\mathbb Q}}),\qquad e=\varepsilon\in M,
 \qquad w=j^*:M\longrightarrow\mathbb Z.
\]
Thus $w(e)=2$. Complex conjugation $c$ acts by
$c(m)=w(m)e-m$. All the assertions below about arbitrary finite-field motives
that use the identification of their band with the Weil-number group retain the
Tate-conjecture hypothesis of R3. The lattice statements, the examples of
abelian varieties, and the calculations with crystals do not assume the Tate
conjecture for higher-codimension cycles.

\paragraph{D.0a (Proved under the Tate identification): audit of the part of R3 used here.}
Let $M$ be the full group of Weil $q$-numbers which are units away from
$p$, modulo roots of unity, and let $M_0=\ker w$. The class
$t=[\sqrt q]$ is Galois invariant, has weight $1$, and satisfies $e=2t$.
Consequently
\[
 M=\mathbb Zt\oplus M_0,\qquad
 M/\mathbb Ze\simeq\mathbb Z/2\oplus M_0.
\]
The weight-zero group $M_0$ is uniquely divisible: every root of a
weight-zero Weil number is again such a number, while the ambiguity
between roots is killed in the quotient by roots of unity. Conjugation
acts as $-1$ on $M_0$. Thus duality of multiplicative-type groups gives
$U\simeq\mu_2\times D(M_0)$, with the second factor an inverse limit of
compact real tori. This verifies the band calculation from the assumed
Tate classification; it is not an unconditional proof of that
classification. The relevant classification and endomorphism statements
are \citet[Proposition 2.6 and Theorem 2.19]{Milne1994}.

At $\ell\ne p$, the \'etale realization with the specified Artin
identification and a Tate basis gives a point of the relative gerbe over
$\mathbb Q_\ell$, so its class there is zero. For the $\mu_2$ component
at $p$ and at infinity, take a supersingular elliptic curve over
$\mathbb F_{q^2}$ with Frobenius $q$. Honda--Tate gives its quaternion
endomorphism algebra with invariants $1/2$ at $p$ and infinity and
zero elsewhere. The same Artin-descent construction written out in
D.3b evaluates the character $t$ of $U$ on this algebra. Since $t$
is precisely the generator of the $\mathbb Z/2$ summand, both local
$\mu_2$ components are nonzero. This checks the local assertion used
from R3; it says nothing about the order of the remaining component.

\paragraph{D.0b (Proved): audit of the real-torus step in R8.}
For a compact real torus $V\simeq
(R^1_{\mathbb C/\mathbb R}\mathbb G_m)^d$, the norm exact sequence gives
\[
 H^1(\mathbb R,V)\simeq(\mathbb Z/2)^d\simeq V(\mathbb R)[2],
 \qquad H^2(\mathbb R,V)=0.
\]
Indeed the relevant terms of the norm sequence are
\[
 \begin{aligned}
 H^1(\mathbb R,R^1\mathbb G_m)
   &=\mathbb R^\times/N\mathbb C^\times=\{\pm1\},\\
 H^1(\mathbb R,\mathbb G_m)&=0,\qquad
 H^2(\mathbb R,R_{\mathbb C/\mathbb R}\mathbb G_m)
   =\operatorname{Br}(\mathbb C)=0 .
 \end{aligned}
\]
Exactness forces $H^2(\mathbb R,R^1\mathbb G_m)=0$.
These identities also appear in
the proof of \citet[Proposition 2.42]{Milne1994}.
For a finite free character lattice with $w(e)=2$, the divisibility
of $e$ can only be $1$ or $2$. Its kernel torus is connected in the
first case. In the second, $e=2t$ with $w(t)=1$ splits the kernel
as a compact torus times $\mu_2$; the nonzero quaternion component
above obstructs a real fibre in the finite-field motivic setting.
In the connected-kernel case, relative real fibre functors form a
torsor under the displayed $H^1$. The parity homomorphisms
$a:M\to\mathbb F_2$ with $a(e)=1$ form a torsor under
$\operatorname{Hom}(M/\mathbb Ze,\mathbb F_2)=V(\mathbb R)[2]$.
This verifies the group calculation behind the fibre/polarization
correspondence used below, with its commutativity and finite-rank
hypotheses explicit.

\paragraph{D.1a (Proved): exact criterion for a rational slope to restrict to $\mu _2$.}
Let $M$ be any subgroup of the group of Weil $q$-numbers, or any subgroup
of its quotient by roots of unity, containing $e=q$. A prime
$\mathfrak p$ over $p$ gives the additive map
\[
 s_{\mathfrak p}:M\longrightarrow\mathbb Q,
 \qquad s_{\mathfrak p}(\lambda)=
 \frac{v_{\mathfrak p}(\lambda)}{v_{\mathfrak p}(q)}.
\]
It satisfies $s_{\mathfrak p}(e)=1$ and
$s_{\mathfrak p}(m)+s_{\mathfrak p}(cm)=w(m)$.
A restriction of this rational cocharacter to $\mu _2$, compatible with the
usual restriction of every integral cocharacter and with clearing odd
 denominators, exists exactly when
\[
 s_{\mathfrak p}(M)\subset\mathbb Z_{(2)}
 =\{a/b\in\mathbb Q:b\text{ odd}\}.
 \tag{D.1}
\]
In that case it is unique and its dual is the homomorphism
$a_{\mathfrak p}:M\to\mathbb F _2$ given by reduction of
$s_{\mathfrak p}$ modulo $2$.

\begin{proof}
For odd $b$, multiplication by $b$ on $\mathbb F _2$ is the identity, so reduction
$a/b\mapsto a\bmod2$ is a well-defined additive map on $\mathbb Z_{(2)}$.
It gives the required homomorphism, with $a_{\mathfrak p}(e)=1$ and
$a_{\mathfrak p}+a_{\mathfrak p}c=w\bmod2$. Conversely suppose
$s_{\mathfrak p}(m)=a/(2^r b)$, where $r\geq1$ and $a,b$ are odd.
On $2^r b m$, slope parity must be $a\bmod2=1$; every additive homomorphism
$M\to\mathbb F _2$ takes that element to zero. This is a contradiction.
Uniqueness follows by multiplication by the odd denominator of each value.
Equivalently, if $D(A)$ denotes the diagonalizable group with character group
$A$, condition (D.1) gives
\[
 \mu _2\longrightarrow D(\mathbb Z_{(2)})
 \longrightarrow D(M)=T_{\mathbb C}.
\]
The first arrow is dual to reduction modulo $2$, and the second to
$s_{\mathfrak p}$; their composite is the precise replacement for the
partially defined notation $i|\mu _2$ on 15:123.
When $M$ has torsion, slope parity vanishes on that torsion and the
map factors through the connected group with character group
$M/M_{\mathrm{tors}}$. Thus the criterion covers subgroups before as
well as after taking the connected quotient.
\end{proof}

\paragraph{D.1b (Proved): the largest simultaneous domain.}
For a Galois-stable group $M$, put
\[
 M^{(2)}=\{m\in M:s_{\sigma\mathfrak p}(m)\in\mathbb Z_{(2)}
       \text{ for every }\sigma\in\operatorname{Gal}(\overline{\mathbb Q}/\mathbb Q)\}.
 \tag{D.2}
\]
This is the largest Galois-stable subgroup on which all the slope restrictions
exist. The set $\{a_{\sigma\mathfrak p}\}$ is a canonical Galois-equivariant
packet on $D(M^{(2)})$. It is independent of the prime initially chosen,
since changing the prime permutes the orbit. This assertion concerns the
quotient group $D(M)\to D(M^{(2)})$, or equivalently the subcategory with that
character group, and does not assert that its fibre functors extend to the
original category.

\begin{proof}
Each inverse image in (D.2) is a subgroup; Galois automorphisms permute the
conditions. D.1a gives existence, uniqueness and the stated equivariance on
this intersection. Every subgroup on which all restrictions exist satisfies
every condition in (D.2), proving maximality. Moreover $e\in M^{(2)}$ cannot
be twice an element of $M^{(2)}$, since its slopes are $1$. Thus the parity
obstruction to a real fibre functor is absent on this subcategory. At each
finite-type torus quotient its compact kernel has vanishing $H^2(\mathbb R,-)$,
so the real gerbe is neutral in the setting of 15:117. The packet assertion is
an assertion about parity homomorphisms, and requires no inverse-limit
neutrality assertion beyond this finite-type statement.
\end{proof}

\paragraph{D.1c (Refuted, with exact repair): the proposed $D(\mathbb Q)$ remedy.}
Passing to the rational-slope protorus does not manufacture the missing
restriction. Indeed
\[
 \operatorname{Hom}(\mu _2,D(\mathbb Q))
       =\operatorname{Hom}(\mathbb Q,\mathbb Z/2)=0.
 \tag{D.3}
\]
Every homomorphism on the right vanishes because $x=2(x/2)$. Thus the canonical
map $D(\mathbb Q)\to\mathbb G_m$ has no lift of the nontrivial element of
$\mathbb G_m[2]$. Rational slopes do define $D(\mathbb Q)\to T$, as in the
Newton homomorphism of \citet[Introduction and \S4]{D_Kottwitz1985}; they do
not define a polarizing $\mu _2$-homomorphism through this source. Lead D1 is
false if it claims that this cover solves the parity problem or that even
slope denominators preclude real fibre functors. The correct restriction
source is $D(\mathbb Z_{(2)})$, on exactly the subcategory (D.2).

\paragraph{D.1d (Proved, computed): the requested $q=25$ example.}
Let $\varpi=2+i$, $q=25$ and
$\lambda=5\varpi^2=15+20i$. In $\mathbb Z[i]$ we have
\[
 (5)=(\varpi)(\overline\varpi),\quad
 v_{\varpi}(\lambda)=3,\quad v_{\overline\varpi}(\lambda)=1,
 \quad\lambda\overline\lambda=625=q^2.
\]
Set $u=\lambda/q=\varpi/\overline\varpi$. The torsion-free group generated by
$q$ and $\lambda$ is $M=\mathbb Ze\oplus\mathbb Zf$, with $e=q$, $f=u$;
its weight map is $(2,0)$ and conjugation is $(e,f)\mapsto(e,-f)$. The
associated $\mathbb Q$-torus is
$\mathbb G_m\times R^1_{\mathbb Q(i)/\mathbb Q}\mathbb G_m$. The two slope
maps are $(1,1/2)$ and $(1,-1/2)$. Therefore
\[
 M^{(2)}=\mathbb Ze\oplus2\mathbb Zf,
 \qquad a_{\mathfrak p}(e)=a_{\mathfrak p}(2f)=1.
\]
The two polarizations of the original real torus are $(a(e),a(f))=(1,0)$ and
$(1,1)$; neither restricts to the displayed slope parity on $M^{(2)}$.
Its real fibre gerbe nevertheless has two isomorphism classes: the kernel
of $\varepsilon$ is the compact one-dimensional torus and
$H^1(\mathbb R,R^1_{\mathbb C/\mathbb R}\mathbb G_m)=\mathbb Z/2$.
All these arithmetic and parity statements are verified by
\texttt{code/verify\_examples.py}, whose output is included. This example also
shows why saying merely that the rational slope is ``not integral'' is too
weak: odd denominators are harmless, whereas this denominator is $2$.

\paragraph{D.1e (Proved): the full category versus a finite subcategory.}
In the full Weil-number group, $e=q$ is twice $\sqrt q$; its class modulo roots
of unity is still twice that class. Hence $a(e)=1$ is impossible on the full
connected group. In the conditional motivic setting of R3 there is no
unmodified real fibre functor on the full category, for any $q$.
On the full even-weight subgroup modulo roots of unity there is, in contrast,
a canonical parity $a=w/2\bmod2$. It is the only parity with $a(e)=1$:
its weight-zero subgroup is divisible, since all roots of a weight-zero
Weil number are again weight-zero Weil numbers, and every homomorphism from
a divisible group to $\mathbb F _2$ vanishes. These statements do not depend
on a choice of slope. They are consistent with the uniqueness statement
for parity-one polarizations on weight-zero motives in
\citet[Lemma 2.46]{Milne1994}. Thus a request for a ``privileged packet''
without specifying slope compatibility has no unique mathematical meaning;
with the compatibility of 15:123 it has the exact obstruction D.1a.

There is nevertheless a useful alternative that works for every $q$
on every finite-type even-weight subcategory. Namely
$a_0=w/2\bmod2$ is defined there, is Galois invariant, and is preserved
by every automorphism respecting $w$ and $e$. The splitting
$M=\mathbb Ze\oplus\ker(w/2)$ makes the kernel of $\varepsilon$ a
compact real torus. By D.0b the real gerbe is neutral, and the R8
correspondence selects the fibre-functor class belonging to $a_0$.
This gives a canonical singleton packet from the even grading alone.
In D.1d it is $(a(e),a(f))=(1,0)$. It is not slope-compatible:
it takes $2f$ to zero, whereas the admissible slope packet takes
$2f$ to one. Thus the existence of an alternative Galois-invariant
selection does not repair the undefined restriction of the original
slope.

\paragraph{D.2a (Proved): complete real-torus orbit calculation.}
Let $T$ be a connected real torus with compact $\ker\varepsilon$, primitive
$e=\varepsilon\in M$, and $w(e)=2$, as above. Its set of polarizing
structures is
\[
 \mathcal P(M,e)=\{a\in\operatorname{Hom}(M,\mathbb F _2):a(e)=1\}.
 \tag{D.4}
\]
If $w(M)=\mathbb Z$, the group
$\operatorname{Aut}(T,\varepsilon,j)$ acts transitively on this set. If
$w(M)=2\mathbb Z$, the orbits are the singleton
$\{w/2\bmod2\}$ and its complement; the complement is empty exactly when
$\dim\ker\varepsilon=0$.

\begin{proof}
The second polarization identity is automatic from $cm=w(m)e-m$ and
$a(e)=1$. All lattice automorphisms preserving $e$ and $w$ commute with $c$,
so they are precisely the real automorphisms in question.
If $w=2w_0$, the identity $w_0(e)=1$ splits
$M=\mathbb Ze\oplus N$, $N=\ker w_0$. Automorphisms are exactly
$1\oplus\operatorname{GL}(N)$. The reduction of
$\operatorname{GL}(N)$ acts transitively on the nonzero elements of
$\operatorname{Hom}(N,\mathbb F _2)$, while fixing zero. This proves the
second assertion, including rank zero.
If $w$ is primitive, choose $x$ with $w(x)=1$ and write
$e=2x+v$, $v\in\ker w$. Primitivity of $e$ implies that some coordinate
of $v$ in a basis of $\ker w$ is odd. Replacing $x$ by $x+k$ changes $v$
by $-2k$, so we can make that coordinate $1$ and hence make $v$ primitive.
Extend $v$ to a basis of $\ker w$, writing
$M=\mathbb Zx\oplus\mathbb Zv\oplus N$. Every $a\in\mathcal P$ has
$a(v)=1$. The automorphisms fixing $x,v$ and sending
$n\mapsto n+\ell(n)v$, with $\ell\in\operatorname{Hom}(N,\mathbb Z)$,
allow arbitrary changes of $a|_N$. The automorphism
$x\mapsto x+v$, $v\mapsto-v$, $n\mapsto n$ changes $a(x)$ by $1$.
All preserve $e,w$. These operations are transitive on (D.4).
\end{proof}

\paragraph{D.2b (Refuted, with an exact finite invariant): rational transitivity.}
Real transitivity does not imply rational transitivity, even in the odd-weight
case. If $T$ is defined over $\mathbb Q$, let $\Delta$ be the finite image
of Galois in $\operatorname{Aut}(M)$. Its rational orbits are precisely the
orbits on (D.4) of the finite group
\[
 \overline{\mathcal A}_{\mathbb Q}
 =\operatorname{im}\bigl(\{A\in\operatorname{GL}(M):Ae=e,\ wA=w,
                    \ A\delta=\delta A\ (\delta\in\Delta)\}
              \longrightarrow\operatorname{GL}(M/2M)\bigr).
 \tag{D.5}
\]
For a possibly disconnected group of multiplicative type with character
module $N$, the corresponding answer uses the image in $\operatorname{GL}(M/2M)$
of Galois-equivariant automorphisms of $N$ preserving $e,w$, where
$M=N/N_{\mathrm{tors}}$. This image can be smaller than (D.5). These formulas
specify the exact finite orbit invariant for a supplied character module;
we do not claim a classification of all such integral Galois modules.

For an explicit example arising from Weil numbers take $q=17$,
\[
 \alpha=1+4i,\qquad \beta=3+2\sqrt{-2},\qquad
 M=\langle\alpha,\overline\alpha,\beta/\overline\beta\rangle.
\]
Both norms are $17$. Independent conjugations of the two quadratic fields
show that the displayed generators are multiplicatively independent:
a relation first forces the exponent of $\beta/\bar\beta$ to vanish,
then the valuations at the two primes over $17$ in $\mathbb Q(i)$ force
the other two exponents to vanish. In the resulting basis $(x,y,t)$,
$e=x+y$, $w=(1,1,0)$; one Galois involution exchanges $x,y$ and fixes $t$,
while another fixes $x,y$ and negates $t$. A rational automorphism must
preserve the three distinct rational Galois eigenspaces. Preservation of
$e$ forces the invariant eigenline to be fixed, and integral invertibility
forces multiplication by $\pm1$ on the other two lines. Thus the only
changes on polarizations are the exchange of $a(x)$ and $a(y)$; $a(t)$
is fixed. The four polarizations form exactly two rational orbits,
\[
 \{(0,1,0),(1,0,0)\},\qquad
 \{(0,1,1),(1,0,1)\}.
\]
The torus is connected, so this also refutes transitivity for
$\operatorname{Aut}(G,\varepsilon,j)$. Honda--Tate realizes the two
weight-one numbers, and the weight-zero ratio belongs to the generated
abelian-motive category. No assertion about higher Tate classes is needed
to verify the lattice example. The even-weight example in D.1d has two
singleton orbits even over $\mathbb R$.

For the full connected Weil-number group of D.1e, the set (D.4) is
empty; for its full even-weight quotient it is a singleton, fixed by
all three automorphism groups. The finite-lattice classification above
concerns finite-type groups, as in the dimension statements on 15:119.
For a general infinite-rank character module the same automorphism
action on $\operatorname{Hom}(M,\mathbb F_2)$ gives the exact orbit
relation, but we do not assert the finite-rank transitivity theorem
without that hypothesis.

\paragraph{D.3a (Proved): inversion and a coherent order-two lift.}
For an abelian band $U$ over a field $k$, in a Galois-cocycle presentation
of a $U$-gerbe of class $h\in H^2(k,U)$, an equivalence inducing inversion
on $U$ exists if and only if $2h=0$. Whenever it exists it can be chosen
to square to the identity, coherently. Thus in this particular inversion
problem there is no additional order-two obstruction after $2h=0$.

\begin{proof}
Inversion carries $h$ to $-h$, so invariance requires $2h=0$.
Write the band additively and represent the gerbe by a continuous cocycle
$c\in Z^2(\Gamma,U(\bar k))$. Its extension has multiplication
\[
 (a,\sigma)(a',\tau)
  =(a+\sigma a'+c(\sigma,\tau),\sigma\tau).
\]
If $2h=0$, choose a continuous $1$-cochain $b$ with $\delta b=-2c$.
The map
\[
 (a,\sigma)\longmapsto(-a+b(\sigma),\sigma)
 \tag{D.6}
\]
is an extension automorphism: substituting into the multiplication law
reduces the assertion exactly to $\delta b=-2c$. Applying it twice gives
$(a,\sigma)$, with no residual cochain. This supplies a strict involution
of the extension presentation and hence a coherent involution of its gerbe.
Changing presentation transports that involution; uniqueness is not claimed.
For a pro-band, the same argument applies to a compatible cocycle
presentation. The obstruction used below is already detected in a single
finite-type pushout, so no inverse-limit exactness is needed for its
nonvanishing.
\end{proof}

\paragraph{D.3b (Proved under the R3 band identification): the fundamental class has infinite order.}
For every prime $p$ and every $q=p^n$, the class of the gerbe of fibre functors
on all finite-field motives, rigidified on Artin motives and the Tate object
as in R3, has infinite order. Consequently its canonical band involution,
which is inversion on $U$, does not lift to any gerbe equivalence, and in
particular not to an involution. The obstruction is already visible using
abelian varieties and their endomorphisms.

\begin{proof}
Fix an integer $r\geq3$ and put
\[
 \alpha=\frac{1+\sqrt{1-4p}}2,\quad K=\mathbb Q(\alpha),\quad
 \pi_r=\alpha^{n(r-1)}\overline\alpha^{\,n}.
 \tag{D.7}
\]
The polynomial $X^2-X+p$ has two distinct roots modulo $p$, and its
discriminant is prime to $p$. Thus $p$ splits in the maximal order of
$K$ and $(\alpha)$ is one of the primes of norm $p$; its conjugate is
the other. We have $\pi_r\bar\pi_r=p^{nr}=q^r$. Hence $\pi_r$ is an
integral Weil $q^r$-number of weight one. Its two slopes are
$(r-1)/r$ and $1/r$. The ratio $\pi_r/\bar\pi_r$ has nonzero valuations,
so no positive power is a root of unity. In particular every
$\mathbb Q(\pi_r^m)$ is the same quadratic field $K$.

By \citet[Th\'eor\`eme 1(i)--(ii), p.~96]{Tate1971}, there is a simple
abelian variety $A_r/\mathbb F_{q^r}$ with this Frobenius. Its division
endomorphism algebra $D_r$ has invariants $\pm1/r$ at the two places over
$p$, with the signs depending on the covariant/contravariant convention,
and zero elsewhere. Therefore its Brauer class has order $r$ and
$\dim A_r=r$. After any extension of degree $m$, the Frobenius is
$\pi_r^m$, the centre remains $K$, and the same invariant formula gives
index $r$ for its simple Honda--Tate factor. Since $A_r$ still has dimension
$r$, it remains simple; dimensions also show that its endomorphism algebra
has not grown. Thus $A_r$ is absolutely simple and $D_r$ is its geometric
endomorphism algebra. This also proves that the use of a finite extension
here preserves the required division-algebra obstruction.

We spell out why this algebra evaluates the actual relative gerbe.
Write $k=\mathbb F_q$, $k'=\mathbb F_{q^r}$ and fix an embedding
$\tau_0:k'\hookrightarrow\overline k$. The Artin fibre functor being
rigidified is
$\omega_{\rm A}(h^0(Z))=H^0(Z_{\overline k},\mathbb Q)$.
For a $\mathbb Q$-algebra $R$, an object of the relative gerbe consists
of a fibre functor $\omega$ with a tensor isomorphism
$\omega|_{\mathcal A}\simeq\omega_{\rm A}\otimes R$ on Artin motives
and a trivialization of the line $\omega(\mathbb Q(-1))$.
Its automorphism band is precisely $U_R$.

Let $B_r=\operatorname{Res}_{k'/k}A_r$ and
$C_r=h^0(\operatorname{Spec}k')$, the finite \'etale algebra object
whose Artin fibre is $\mathbb Q^{\operatorname{Hom}_k(k',\overline k)}$.
There is a canonical action
$C_r\otimes h^1(B_r)\to h^1(B_r)$ in the rational motive category
over $k$. To check descent, base change to $k'$. The abelian variety
$B_{r,k'}$ is the product of the $r$ Galois conjugates of $A_r$.
The $r$ factors of $C_{r,k'}$ act by the orthogonal projection
idempotents on the corresponding $h^1$ summands. Galois permutes
these factors and projections together, so the action is invariant.
It descends as a rational correspondence: pullback followed by
transfer is multiplication by $r$, and averaging the invariant
correspondence divides by this invertible integer. The unit and
associativity identities can be checked after pullback, which is
injective by the same transfer identity.

Every $d\in D_r=\operatorname{End}_{k'}^0(A_r)$ gives an endomorphism
$\operatorname{Res}_{k'/k}(d)$ of $B_r$ over $k$; after base change
these are the diagonal conjugates of $d$. This action commutes with
$C_r$. For a relative fibre functor $(\omega,\iota,t)$, the specified
Artin isomorphism $\iota$ therefore supplies a distinguished
idempotent $e_{\tau_0}$ of $\omega(C_r)$, and
\[
 V_\omega=e_{\tau_0}\omega(h^1(B_r))
\]
is a vector bundle of rank $2r$ over $R$, with a left action of
$D_r^{\rm op}$, because pullback on our contravariant $h^1$ reverses
composition of endomorphisms.
This construction commutes with every scalar extension $R\to R'$
and with every isomorphism of relative fibre functors.

Now pull the gerbe back from $\mathbb Q$ to the coefficient field $K$.
For a $K$-algebra $R$, the two idempotents of
$K\otimes_{\mathbb Q}R\simeq R\times R$ separate the operator action
of the centre $K$ on $V_\omega$. Take the idempotent for the identity
embedding $K\hookrightarrow R$, and call its image $V_{\omega,{\rm id}}$.
It has rank $r$: the Frobenius polynomial of $A_r$ is the quadratic
minimal polynomial of $\pi_r$ to the $r$th power, so its two
$K$-embedding components have equal rank $r$. The action
\[
 D_r^{\rm op}\otimes_K R\longrightarrow
       \operatorname{End}_R(V_{\omega,{\rm id}})
\]
is an isomorphism. This can be checked fpqc locally, where a unital
representation of the matrix algebra $M_r$ on a rank-$r$ module is
the standard representation and both algebras have rank $r^2$.

Let $\mathscr S(D_r^{\rm op})$ be the gerbe over $K$ whose objects over $R$ are
rank-$r$ vector bundles with such a splitting action and whose arrows
are $D_r^{\rm op}$-linear isomorphisms. It is a $\mathbb G_{m,K}$-gerbe:
the commutant of the full matrix algebra is exactly the scalars.
Fix the Brauer-to-gerbe convention in which the gerbe of left
splitting modules of an algebra $D$ has class $[D]$; its class is
then $[D_r^{\rm op}]=-[D_r]$. The construction
$\omega\mapsto V_{\omega,{\rm id}}$ is now an actual morphism of
gerbes over $K$, not merely a comparison of endomorphism algebras.
An automorphism in $U_R$ commutes with the $D_r^{\rm op}$-action and hence
acts by a scalar. These scalar actions give the morphism of bands
\[
 \chi_r:U_K\longrightarrow\mathbb G_{m,K}.
\]
In the Weil-number character description this descent is also visible
directly: choose $\lambda_r$ with $\lambda_r^r=\pi_r$. It is a Weil
$q$-number. Its class modulo roots of unity is fixed by
$\operatorname{Gal}(\overline{\mathbb Q}/K)$, since every conjugate over
$K$ differs from it by an $r$th root of unity. This is the central
character on the chosen Artin summand, restricted further to $U$.
Thus the scalar character above has the required field of definition.
Naturality of the class of a banded gerbe gives the pushout identity
\[
 (\chi_r)_*(h|_K)=-[D_r]\in H^2(K,\mathbb G_m)=\operatorname{Br}(K).
 \tag{D.8}
\]
The Tate trivialization was retained in the source throughout; the
construction forgets it and thus defines the same functor after that
rigidification. Neither a change of the geometric base field nor
neutrality of an unrelated Tannakian category has been substituted
for this morphism. A hypothetical integer annihilating $h$ must
annihilate every (D.8), and hence must be divisible by every $r\geq3$.
No nonzero integer has this property. Finally the character identity
$cm=w(m)e-m$ shows that the canonical involution is inversion after
quotienting by $e$; D.3a now excludes a lift.
\end{proof}

The arithmetic of (D.7) was checked exactly for $p=2,3,5,7,11$ and
$r=3,4,5,6,7$; this finite check supports the displayed calculations and
is not offered as a proof for all primes. The general proof above uses
only $X^2-X+p$ and the local invariant theorem. In particular $p=2,3$
are included without an exception. The $\mu _2$-projection can still have
order $2$; it does not control the order of the compact-torus component.
The prior reading's suggestion that the reason on 15:125 is erroneous
because the class should have order dividing $2$ is therefore withdrawn.
The transcribed French phrase ``qui serait ... d'ordre divisant $2$'' can
correctly express a necessary consequence of a proposed lift. We do not
alter the transcription to obtain this interpretation.

\paragraph{D.4a (Proved under explicit compatibility hypotheses): what the distribution condition says.}
Fix a polarizable Tannakian category over a base $S$, an even grading,
a polarization $\Pi$ on it, and specialization tensor functors to the
closed-point categories. Suppose these functors preserve $\Pi$. Suppose
also that the slope parity maps defining $v(s')$ exist as in D.1a and
that the real fibre functor selected at each $s'$ is the one corresponding
to the specialized polarization. If the fibre-functor torsor classes
$\xi(s')$ agree, the polarizations transported by $v(s')$ agree.

\begin{proof}
Choose a representative real fibre functor for the common torsor class.
An isomorphism to that representative transports the same system of
positive forms $\Pi$ on every object. A different isomorphism differs
by a tensor automorphism and gives an isomorphic polarization. Thus both
specializations yield the same polarization on the twisted real group.
The assertion is a compatibility statement about specialization and
positive forms; constancy of $\xi$ alone does not supply these hypotheses.
\end{proof}

As a universal statement using the unqualified $i|\mu _2$ of 15:183--185,
the construction is not defined: D.1d is an even-weight finite-field
example where that arrow does not exist. We have not proved the general
specialization/positivity assertion for all geometric motives over every
finite-type $\mathbb Z$-scheme. The remaining gate is the existence and
specialization compatibility of the algebraic polarizations invoked on
15:185, after restricting to the admissible domain (D.2) or specifying
another selection rule. A finite parity calculation does not settle that
geometric assertion. \citet[Remark 2.47]{Milne1994} explicitly distinguishes
abstract polarizability from the algebraic construction supplied by the
standard conjectures.

\paragraph{D.4b (Proved as an equivalence; finite-field conclusion conditional): the effectivity claim.}
For a fixed tensor category and its specialization maps $u(s')$, define the
spectral cone to consist of characters having nonnegative pairing with
every conjugate of every normalized valuation of the specialized Frobenius.
An algebraic number which is a unit away from its residue characteristic
is integral exactly when these valuations are nonnegative. Therefore
recovery of the actual effective cone from those maps is exactly the
assertion that simultaneous integrality of all specialized eigenvalues
implies effectivity; the reverse implication is the usual integrality
property of the effective realizations. The statement is an equivalence
of two descriptions of the same missing implication, not a proof of that
implication.

Over $\mathbb F_q$, assume the Tate classification of simple pure motives
and semisimplicity. If a simple motive has integral Weil eigenvalue
$\lambda$ of weight $i\geq1$, choose an $i$th root $\alpha$ of $\lambda$.
Every conjugate of $\alpha$ has absolute value $q^{1/2}$ and $\alpha$ is
integral. Honda--Tate gives $A/\mathbb F_q$ with $\alpha$ on $H^1(A)$;
then $\lambda$ occurs on $H^i(A^i)$ by K\"unneth. A simple factor of this
effective motive has the same Weil number, so Tate classification
identifies it with the given simple motive. Weight zero is handled by a
finite \'etale scheme, since an integral weight-zero Weil number is a root
of unity. An integral nonzero number cannot have negative weight, because
its nonzero integral norm would have absolute value less than $1$.
Taking sums proves the finite-field assertion under the stated hypotheses.
Without full faithfulness/classification, equality of Frobenius spectra
alone does not identify the given motive with that effective factor.

\paragraph{D.4c (Proved): a test for specialization to preserve a polarization.}
Let $(\mathcal C,w,\mathbb T)$ be a real Tate triple with polarization
$\Pi$, and let $F:\mathcal C\to\mathcal Q$ be a quotient by a closed
subgroup $H\subset\pi(\mathcal C)$ containing the weight torus.
Assume that $\mathcal Q$ is semisimple. A sufficient condition for a
polarization $\Pi_{\mathcal Q}$ with $F\Pi=\Pi_{\mathcal Q}$ is the following
concrete positivity test:
\[
 \text{for every }Y\in\mathcal C^H
 \text{ and }\psi\in\Pi(Y),\qquad
 F(\psi)\text{ is positive definite on }F(Y).
 \tag{D.9}
\]
Here $F(Y)$ is regarded as an ordinary real vector space because the
fundamental group of $\mathcal Q$ acts trivially on it. Moreover, for any
tensor functor to a semisimple target in which every object is a direct
summand of an image, preservation of a polarization is equivalent to
every image $F(\psi)$ being a Weil form. This gives an actual test, rather
than assuming the desired compatibility in different notation.

\begin{proof}
For $X$ pure of weight $m$ and $\psi:X\otimes X\to\mathbb T^{-m}$
in $\Pi(X)$, form the weight-zero object
$E=\underline{\operatorname{End}}(X)$ and its trace form
\[
 T^\psi(a,b)=\operatorname{Tr}(a^\psi b).
\]
The restriction of this form to $E^H$ belongs to $\Pi(E^H)$.
The quotient property identifies
$\operatorname{End}_{\mathcal Q}(F(X))$ with the real vector space
$F(E^H)$, and tensor functoriality identifies its trace form with
$F(T^\psi|_{E^H})$. Condition (D.9) therefore makes $F(\psi)$ a Weil
form. Every target object is a direct summand of some $F(X)$, so these
forms and their restrictions define a target polarization. To check
independence of the presentation, apply the same construction to the
direct sum of two presentations: the two resulting restricted forms
belong to the same positivity class there. Conversely, a functor
preserving a polarization takes positive Weil forms to positive Weil
forms by definition. This is the criterion and construction of
\citet[Lemma 1.2 and Proposition 1.4, pp.~602--603]{D_Milne2002};
the argument shows precisely why additional endomorphisms after
specialization are the part on which positivity must be checked.
\end{proof}

\paragraph{D.4d (Known with derivation, and with the stated 2026 input): geometric cases of the distribution condition.}
Fix a prime of $\overline{\mathbb Q}$ over $p$ and write
$\mathbb F=\overline{\mathbb F}_p$. On the categories of Lefschetz
motives, with rational correspondences generated by divisor classes
modulo numerical equivalence, the reduction functor
\[
 R_{\rm L}:\operatorname{LCM}(\overline{\mathbb Q})
        \longrightarrow\operatorname{LMot}(\mathbb F)
\]
preserves the canonical polarizations. For motives with all algebraic
correspondences, the same assertion holds for reduction of CM abelian
motives if the rational Hodge conjecture holds for every complex CM
abelian variety. Thus D.4a applies to these reductions wherever the
slope parity used to select the real fibres is defined. In particular,
it applies to the even subcategories generated by a CM abelian scheme
over an open subscheme of a number ring, with its extending
correspondences and relatively ample polarizations, at its admissible
closed points. The equality concerns the polarization structures after
transport to a common real fibre; representing homomorphisms have the
conjugacy ambiguity already stated on 15:183.

\begin{proof}
The first assertion is exactly the reduction assertion in
\citet[Remark 3.7, pp.~608--609]{D_Milne2002}. Its construction uses the
Rosati trace form on the additional endomorphisms; positivity of that
form verifies (D.9). The geometric primitive pairings are then positive
on Lefschetz classes, and tensor operations and direct summands give
the polarization on the whole Lefschetz category. For all algebraic
correspondences, \citet[Theorem 2.1, p.~604]{D_Milne2002} states that,
under the specified Hodge hypothesis, reduction takes the canonical
CM polarization to the canonical finite-field polarization and makes
every geometric Weil form positive. These are precisely the
preservation hypotheses of D.4a, with source and target now specified.

For a CM abelian scheme with good reduction, restriction to a geometric
closed point is this reduction functor on the chosen objects and
correspondences. Relatively ample divisors specialize to ample divisors,
and cup product, trace, and the associated forms commute with smooth
proper specialization. Passing from a geometric closed point to its
finite residue field takes the corresponding endomorphism spaces to
subspaces of those over $\mathbb F$; restriction of a positive definite
trace form remains positive definite. Thus the same preservation
holds at the finite closed point. Applying D.4a gives the asserted
coincidence when $\xi(s')$ is fixed.
\end{proof}

For the date of this research, the Hodge hypothesis in the second case
must be described using the current literature. Theorem I of
\citet{D_VanHoftenSempliner2026}, a primary preprint whose statement and
proof chain we opened, asserts it for every complex CM abelian variety.
Its coefficient field for cycle classes is $\mathbb Q$; it does not
assert an integral Hodge conjecture. Invoking that preprint theorem
removes the Hodge hypothesis in this CM/abelian case. We explicitly
retain it as a cited 2026 input: this run checked its applicability and
the elementary surface-to-cycle reduction, but did not independently
reprove every analytic theorem used in its theta-integral argument.
Theorem 3.3 and Corollary 3.4 of \citet[pp.~607--608]{D_Milne2002}
then give the geometric polarization and the Hodge standard
conjecture for abelian motives over algebraically closed fields.
None of these results states the specialization assertion for arbitrary
geometric motives on every finite-type $\mathbb Z$-scheme.

\paragraph{D.4e (Proved under the cited 2026 theorem): finite-field effectivity for motives of abelian type.}
Let $\mathcal A_q$ be the semisimple category of pure numerical motives
over $\mathbb F_q$, with rational coefficients, generated by abelian
varieties, Artin motives and Tate objects. Assume Theorem I of
\citet{D_VanHoftenSempliner2026}. An object of $\mathcal A_q$ is effective
if and only if all of its Frobenius eigenvalues are algebraic integers.

\begin{proof}
\citet[Theorem 7.1, author version p.~25]{D_Milne1999} derives the Tate
conjecture for all abelian varieties over finite fields from exactly the
Hodge assertion supplied by that Theorem I. We use this theorem with
the cofinal-field correction recorded on the author's errata page:
the auxiliary CM fields can be enlarged to contain a totally real
field of degree greater than two in which $p$ is inert. The unrestricted
Hodge hypothesis is unchanged by this enlargement.

Every product of abelian varieties is an abelian variety. Hence the
Tate assertion applies to the product of any two varieties used in
constructing objects of $\mathcal A_q$. Its cycle statement, followed
by the relevant algebraic projectors, identifies morphisms of numerical
motives with Frobenius-equivariant maps of their $\ell$-adic
realizations, for $\ell\ne p$. The same conclusion with Artin factors
follows after a finite extension, and descends: average a correspondence
over the finite Galois group and divide by the extension degree,
which is allowed for rational coefficients. Tate twists merely change
the cohomological degree and the twist in this comparison. Thus the
full-faithfulness/classification input needed in D.4b holds on
$\mathcal A_q$.

For clarity, one can use full faithfulness directly. Let $M$ be a
simple constituent and $\lambda$ one of its eigenvalues. The
construction in D.4b gives an effective abelian--Artin motive $N$
with eigenvalue $\lambda$. The polynomial of Frobenius on $M$ is a
power of the irreducible polynomial of $\lambda$: its Frobenius
endomorphism is central, and a simple object has a field as the centre
of its endomorphism algebra. The $\ell$-adic Frobenius modules of
$M$ and $N$ consequently have a common simple constituent and a
nonzero equivariant map. Full faithfulness gives a nonzero morphism
of motives. Semisimplicity makes $M$ a direct summand of $N$.
The cases of zero and negative weight, and finite direct sums, were
proved in D.4b. Conversely, the eigenvalues of every effective
summand of smooth projective cohomology are integral.
\end{proof}

The conclusion of D.4e is for the specified category $\mathcal A_q$.
For a general numerical motive over $\mathbb F_q$, the full
faithfulness used in its proof has not been supplied here. The
remaining finite-field gate in D.4b is therefore restricted to
motives outside this category; the 2026 abelian result is not a
proof of the Tate conjecture for arbitrary varieties.

\paragraph{D.4f (Proved under the stated standard framework): the all-motives effectivity criterion detects the Tate conjecture.}
The remaining classification input in D.4b can be characterized
precisely. Fix $q$ and suppose that the category $\mathcal N_q$ of
pure numerical motives over $\mathbb F_q$ has the usual algebraic
Kunneth decomposition, is semisimple with finite-dimensional
rational morphism spaces, and has a faithful $\ell$-adic
realization factoring through numerical equivalence, for some
$\ell\ne p$. The last condition is essential: without it the
phrase \emph{eigenvalues of a numerical motive} has not been
justified by this realization.
Use the usual effective subcategory of direct summands of finite
sums of motives of smooth projective varieties. In this framework
the following are equivalent:
\begin{enumerate}
\item every motive whose Frobenius eigenvalues are algebraic
integers is effective;
\item every motive on which Frobenius is the identity is effective;
\item for every smooth projective $X/\mathbb F_q$ and every $r$,
the cycle map is onto the Frobenius-fixed subspace of
$H^{2r}(X_{\overline{\mathbb F}_q},\mathbb Q_\ell(r))$.
\end{enumerate}
Here (3) is the cycle-class form of the Tate conjecture on all
varieties over the fixed field; the framework already supplies
semisimplicity of the Frobenius actions used below.

\begin{proof}
The implication (1)$\Rightarrow$(2) is immediate. To prove
(2)$\Rightarrow$(3), first note that Frobenius is central in the
endomorphism algebra of every object: its realization commutes
with every correspondence, and the realization is faithful.
The centre of a finite-dimensional semisimple rational algebra
is a product of number fields. Frobenius consequently has a
square-free minimal polynomial. In particular, its fixed part
is cut out by a rational polynomial projector.

More explicitly, put $N=h^{2r}(X)(r)$. If its Frobenius minimal
polynomial has a factor $T-1$, write it as $(T-1)g(T)$ with
$g(1)\ne0$ and choose
\[
 P(T)\equiv1\pmod{T-1},\qquad P(T)\equiv0\pmod{g(T)}.
\]
The Chinese remainder theorem supplies $P\in\mathbb Q[T]$.
Let $M=\operatorname{im}P(F_N)$; if $T-1$ is absent, put $M=0$.
Then $F_M=1$, and its realization is exactly the fixed subspace
in (3). By (2), $M$ is effective.

Every effective pure motive of weight zero is Artin. Indeed
it is a summand of a finite sum of $h(Y)$ with no Tate
twists. The weight decomposition removes every $h^j(Y)$
with $j>0$, leaving only $h^0(Y)$; this last motive is that
of the finite etale scheme of geometric connected components.
For an Artin motive, a Frobenius element topologically
generates the Galois group of $\mathbb F_q$. Hence $F_M=1$
implies $M\simeq\mathbf1^{\oplus d}$.
The inclusions of these unit summands into $h^{2r}(X)(r)$
are rational algebraic cycles modulo numerical equivalence.
Their realizations form a basis of the fixed subspace.
This proves (3), including surjectivity after tensoring the
rational cycle space with $\mathbb Q_\ell$.

Conversely assume (3). Applied to products, it identifies
correspondences, followed by the given Kunneth projectors,
with all Frobenius-equivariant maps between realizations:
\[
 \operatorname{Hom}_{\mathcal N_q}(M,N)\otimes\mathbb Q_\ell
  \simeq
 \operatorname{Hom}_{F}(R_\ell M,R_\ell N).
\]
Injectivity after scalar extension follows from the
nondegenerate numerical trace pairing. Namely, choose
rational dual bases of the two correspondence spaces
$\operatorname{Hom}(M,N)$ and $\operatorname{Hom}(N,M)$
for composition followed by trace. A purported
$\mathbb Q_\ell$-linear relation between the realized
first basis, paired with the realized second basis,
has every coefficient zero. The realization preserves
these rational traces.
Surjectivity is exactly the Tate cycle statement on the
product in the appropriate degree and twist. Thus the
full-faithfulness input in D.4b holds. For an integral
eigenvalue of weight $i>0$, Honda--Tate applied to an
$i$th root supplies an effective abelian motive with that
eigenvalue. A common Frobenius constituent gives a nonzero
map from each simple constituent of the proposed motive
to an effective motive. The displayed isomorphism gives
such a map over $\mathbb Q$ as well: a nonzero tensor
product of a rational morphism space cannot arise from
the zero space. Semisimplicity splits it. Weight zero
and the exclusion of negative weight are as in D.4b.
This proves (1).
\end{proof}

The argument sharpens the remaining issue: even the
Frobenius-identity weight-zero case of the effectivity
claim would supply every Tate class in this framework.
The known abelian-type theorem D.4e cannot be extended
to all numerical motives merely by noting that their
eigenvalues are realizable on abelian varieties.

\paragraph{D.4g (Proved under the cited 2026 theorem): a common geometric polarization for arbitrary abelian families.}
The geometric preservation case can be enlarged beyond CM
families. Let $S$ be integral and of finite type over
$\mathbb Z$, and let $A_1,\ldots,A_t$ be abelian schemes
over $S$ equipped with relatively ample polarizations.
Let $\mathcal C$ be a pure Tate category generated under
direct sums, tensor operations and direct summands by
their degree-one motives and the Tate object, using
relative algebraic correspondences and projectors for
which exact faithful specialization tensor functors
\[
 F_s:\mathcal C\longrightarrow
          \operatorname{Mot}(\overline{k(s)};\mathrm{Ab})
\]
are defined and preserve the weight grading, the Tate
object and its fixed Tate structure. Thus cycles and projectors are assumed to
extend in the given family; a generic correspondence
with no specified specialization is not being used.
Assume Theorem I of \citet{D_VanHoftenSempliner2026}.
There is a common polarization $\Pi_{\mathcal C}$ such
that every $F_s$ preserves it. It is the unique
polarization containing the specified geometric forms
on the generating degree-one motives and the Tate
object. On the even subcategory, the distribution
assertion D.4a therefore holds with these geometric
polarizations and the fibre/polarization identification
specified there, whenever the chosen slope parities
exist.

\begin{proof}
By \citet[Theorem 3.3 and Corollary 3.4,
pp.~607--608]{D_Milne2002}, the CM Hodge hypothesis
supplied by the cited 2026 theorem gives the unique
geometric polarization $\Pi_s$ on the full category
of abelian motives over each algebraically closed
residue field. It also identifies numerical and
$\ell$-adic homological equivalence there. Theorem 2.1
ensures positivity of every geometric degree-one form.
The relatively ample forms in our statement specialize
to exactly these forms.

We give the tensor argument that ensures independence
of $s$. More generally, suppose an exact faithful tensor
functor $F:\mathcal C\to(\mathcal D,\Pi_{\mathcal D})$
carries fixed forms on a set of tensor generators to
positive forms. Work in each pure weight.
Tensor products, duals and direct sums
give a form on any expression in the generators.
If $X$ is a direct summand of that expression, restrict
the form to $X$. Its image under $F$ is a positive,
nondegenerate form, so the restricted form in
$\mathcal C$ is nondegenerate: faithfulness and
exactness reflect the invertibility of its associated
map to the dual. Its trace form on
$\operatorname{End}_{\mathcal C}(X)$ is positive
because this endomorphism space embeds in
$\operatorname{End}_{\mathcal D}(F X)$.

Thus every object has a Weil form whose image is
positive. Define the positive class on $X$ to be the
forms whose image belongs to $\Pi_{\mathcal D}(F X)$.
These classes are nonempty and satisfy the sum and
tensor compatibility axioms because their images do.
They form a polarization on $\mathcal C$; equivalently
this is the pullback construction implicit in
\citet[Lemma 1.2 and its preceding definition,
pp.~602--603]{D_Milne2002}.
Its classes are determined by the forms on the
generators, since the construction by sums, tensors,
duals and summands reaches every object.

Apply this argument to every $F_s$. In each case the
same relative generator forms and the same restricted
forms constructed from them are positive. Two
polarizations containing a common Weil form on every
object have the same compatibility classes. Hence the
pullbacks of all $\Pi_s$ coincide; call their common
value $\Pi_{\mathcal C}$.

The needed faithfulness of a specialization on this
specified category can also be checked geometrically.
A relative correspondence acts on the smooth proper
$\ell$-adic local systems. Vanishing at one geometric
point of connected $S$ means vanishing everywhere.
Numerical equivalence is homological equivalence on
the abelian fibres by Theorem 3.3, so quotienting by
that relation makes each specialization faithful.
For any two points of $S$, choose $\ell$ different
from their residue characteristics. Both lie in
the integral, hence connected, open $S[1/\ell]$,
where the preceding argument applies; the numerical
relation is independent of that choice.
The existence and typing
of the specified specialization functors remain part
of the statement, as they must for motives over a
base.

Finally restrict to even weights and apply D.4a.
After choosing an isomorphism to a representative
of the common real fibre class $\xi$, the same
polarization $\Pi_{\mathcal C}$ is transported at
every point. Changes of the chosen isomorphism give
the prescribed real inner-conjugacy ambiguity.
\end{proof}

This is a result for abelian-type families and their
extending algebraic correspondences. It does not
deduce from constancy of an unstructured torsor
class that an arbitrary choice of $\mu_2$-map is
geometrically positive. The requested assertion for
general geometric motives still needs precisely
that positivity and specialization compatibility,
in addition to a defined slope-parity map.
